\documentclass[aspectratio=169]{beamer}
\usetheme{Madrid}
\usecolortheme{whale}
\usepackage[utf8]{inputenc}
\usepackage[spanish]{babel}
\usepackage{amsmath,amsfonts,amssymb}
\usepackage{graphicx}
\usepackage{booktabs}
\usepackage{listings}
\usepackage{xcolor}

\lstdefinestyle{rcode}{
  language=R,
  basicstyle=\ttfamily\footnotesize,
  backgroundcolor=\color{gray!10},
  frame=single,
  rulecolor=\color{gray!50},
  breaklines=true,
  keywordstyle=\color{blue!80!black}\bfseries,
  commentstyle=\color{green!50!black}\itshape,
  stringstyle=\color{red!70!black},
  showstringspaces=false,
  numbers=none,
  tabsize=2
}

\title{Métodos de Análisis de Datos 1}
\subtitle{Clase 13: Visualización con R}
\author{Depto.\ de Matemática -- UNS}
\institute{Universidad Nacional del Sur}
\date{1er Cuatrimestre 2026}

\AtBeginSection[]{
  \begin{frame}{Contenidos}
    \tableofcontents[currentsection]
  \end{frame}
}

\begin{document}

\begin{frame}
  \titlepage
\end{frame}

\begin{frame}{Contenidos}
  \tableofcontents
\end{frame}

%=============================================================
\section{Sistemas de graficado en R}
%=============================================================

\begin{frame}{Los tres sistemas principales}
  \begin{center}
    \begin{tabular}{lll}
      \toprule
      \textbf{Sistema} & \textbf{Paradigma} & \textbf{Características} \\
      \midrule
      \texttt{base R} & Tradicional & Simple, funciones como \texttt{plot()}, \texttt{hist()} \\
      \texttt{lattice} & Fórmulas & Paneles múltiples, datos condicionados \\
      \texttt{ggplot2} & Gramática & Capas declarativas, muy flexible y elegante \\
      \bottomrule
    \end{tabular}
  \end{center}
  \bigskip
  \begin{block}{Enfoque del curso}
    Usaremos \texttt{ggplot2} como sistema principal por su coherencia conceptual y la calidad de sus gráficos. También veremos funciones de \texttt{base R} para análisis rápido.
  \end{block}
\end{frame}

%=============================================================
\section{Gramática de los Gráficos}
%=============================================================

\begin{frame}{Gramática de los Gráficos (Wilkinson, 2005)}
  \begin{block}{Idea central}
    Un gráfico es la composición de capas independientes. Cada capa describe un aspecto del gráfico.
  \end{block}
  \bigskip
  \textbf{Componentes de la gramática en \texttt{ggplot2}:}
  \begin{enumerate}
    \item \textbf{Data}: el DataFrame con los datos.
    \item \textbf{Aesthetics} (\texttt{aes}): mapeo de variables a propiedades visuales ($x$, $y$, color, tamaño, forma).
    \item \textbf{Geometries} (\texttt{geom\_*}): tipo de gráfico (puntos, líneas, barras, boxplot).
    \item \textbf{Statistics} (\texttt{stat\_*}): transformaciones estadísticas.
    \item \textbf{Scales}: control de ejes, colores, tamaños.
    \item \textbf{Facets}: paneles múltiples por grupo.
    \item \textbf{Theme}: apariencia general.
  \end{enumerate}
\end{frame}

\begin{frame}[fragile]{Estructura básica de un gráfico ggplot2}
\begin{lstlisting}[style=rcode]
library(ggplot2)

# Estructura general
ggplot(data = mis_datos,
       aes(x = variable_x, y = variable_y)) +
  geom_TIPO() +
  labs(title = "Titulo", x = "Eje X", y = "Eje Y") +
  theme_bw()

# Ejemplo: scatter plot
ggplot(data = iris,
       aes(x = Sepal.Length, y = Sepal.Width,
           color = Species)) +
  geom_point(size = 2, alpha = 0.7) +
  labs(title = "Iris: longitud vs. ancho del sepalo",
       x = "Longitud (cm)", y = "Ancho (cm)") +
  theme_bw()
\end{lstlisting}
\end{frame}

%=============================================================
\section{Gráficos univariados}
%=============================================================

\begin{frame}[fragile]{Histograma y densidad con ggplot2}
\begin{lstlisting}[style=rcode]
library(ggplot2)

datos <- data.frame(x = rnorm(500, mean = 50, sd = 10))

# Histograma con curva de densidad superpuesta
ggplot(datos, aes(x = x)) +
  geom_histogram(aes(y = after_stat(density)),
                 bins = 25, fill = "steelblue",
                 color = "white") +
  geom_density(color = "red", linewidth = 1) +
  labs(title = "Histograma con densidad",
       x = "Valor", y = "Densidad") +
  theme_bw()
\end{lstlisting}

\medskip
\textbf{Nota}: \texttt{after\_stat(density)} escala el histograma a densidad para que sea comparable con la curva KDE (Kernel Density Estimate).
\end{frame}

\begin{frame}[fragile]{Boxplot y violin plot}
\begin{lstlisting}[style=rcode]
library(ggplot2)

# Boxplot por grupo
ggplot(iris, aes(x = Species, y = Sepal.Length,
                 fill = Species)) +
  geom_boxplot(alpha = 0.7, outlier.colour = "red") +
  labs(title = "Longitud del sepalo por especie",
       x = "Especie", y = "Longitud (cm)") +
  theme_bw() +
  theme(legend.position = "none")

# Violin plot (reemplazar geom_boxplot por):
# geom_violin(alpha = 0.7, trim = FALSE) +
# geom_boxplot(width = 0.1, fill = "white")
\end{lstlisting}
\end{frame}

%=============================================================
\section{Gráficos bivariados}
%=============================================================

\begin{frame}[fragile]{Scatter plot con regresión}
\begin{lstlisting}[style=rcode]
library(ggplot2)

ggplot(mtcars, aes(x = wt, y = mpg)) +
  geom_point(color = "steelblue", size = 3,
             alpha = 0.7) +
  geom_smooth(method = "lm", se = TRUE,
              color = "red", fill = "red",
              alpha = 0.15) +
  labs(title = "Consumo vs. peso del vehículo",
       x = "Peso (1000 lbs)",
       y = "Millas por galón (mpg)") +
  theme_bw()
\end{lstlisting}

\medskip
\texttt{se = TRUE} muestra el intervalo de confianza del 95\% para la línea de regresión.
\end{frame}

\begin{frame}[fragile]{Heatmap de correlaciones}
\begin{lstlisting}[style=rcode]
library(ggplot2)
library(reshape2)

# Calcular matriz de correlacion
corr_mat <- cor(mtcars)

# Pasar a formato largo
corr_long <- melt(corr_mat)

# Heatmap
ggplot(corr_long, aes(x = Var1, y = Var2,
                      fill = value)) +
  geom_tile() +
  scale_fill_gradient2(low = "blue", mid = "white",
                       high = "red", midpoint = 0) +
  labs(title = "Matriz de correlaciones - mtcars",
       x = "", y = "", fill = "r") +
  theme_bw() +
  theme(axis.text.x = element_text(angle = 45,
                                   hjust = 1))
\end{lstlisting}
\end{frame}

%=============================================================
\section{Paneles y temas}
%=============================================================

\begin{frame}[fragile]{Facetas: paneles múltiples}
\begin{lstlisting}[style=rcode]
library(ggplot2)

# facet_wrap: una variable de agrupacion
ggplot(iris,
       aes(x = Sepal.Length, y = Petal.Length)) +
  geom_point(aes(color = Species), alpha = 0.6) +
  geom_smooth(method = "lm", se = FALSE) +
  facet_wrap(~ Species, ncol = 3) +
  labs(title = "Relacion longitud sepalo-petalo",
       x = "Longitud sepalo (cm)",
       y = "Longitud petalo (cm)") +
  theme_bw() +
  theme(legend.position = "none")
\end{lstlisting}
\end{frame}

\begin{frame}{Temas (\texttt{theme}) en ggplot2}
  Los temas controlan la apariencia general del gráfico.
  \bigskip
  \textbf{Temas predefinidos más usados:}
  \begin{itemize}
    \item \texttt{theme\_grey()} — predeterminado (fondo gris).
    \item \texttt{theme\_bw()} — fondo blanco, cuadrícula gris. Recomendado para publicaciones.
    \item \texttt{theme\_classic()} — sin cuadrícula, ejes limpios.
    \item \texttt{theme\_minimal()} — minimalista, sin bordes.
    \item \texttt{theme\_void()} — sin ejes ni fondo (útil para mapas).
  \end{itemize}
  \bigskip
  \begin{block}{Personalización}
    Cualquier elemento se puede ajustar con \texttt{theme(element = element\_*(...))}. Por ejemplo: tamaño de texto, posición de leyenda, ángulo de etiquetas.
  \end{block}
\end{frame}

%=============================================================
\section{Exportación}
%=============================================================

\begin{frame}[fragile]{Guardar gráficos}
\begin{lstlisting}[style=rcode]
# Con ggsave (recomendado para ggplot2)
p <- ggplot(iris, aes(x = Sepal.Length,
                      y = Sepal.Width)) +
  geom_point()

ggsave(filename = "mi_grafico.pdf",
       plot = p,
       width = 8, height = 5,
       units = "in")

ggsave("mi_grafico.png", plot = p,
       width = 8, height = 5, dpi = 300)

# Con funciones base R
pdf("grafico_base.pdf", width = 8, height = 5)
hist(iris$Sepal.Length, main = "Histograma",
     xlab = "Longitud (cm)")
dev.off()
\end{lstlisting}
\end{frame}

%=============================================================
\section{Resumen}
%=============================================================

\begin{frame}{Resumen de la clase y de la Unidad 2}
  \begin{block}{Visualización con R -- conceptos clave}
    \begin{itemize}
      \item Gramática de los Gráficos: datos + estéticas + geometrías + temas.
      \item Construcción por capas con el operador \texttt{+}.
      \item Facetas para paneles múltiples: \texttt{facet\_wrap()}, \texttt{facet\_grid()}.
      \item Exportación con \texttt{ggsave()} en PDF o PNG.
    \end{itemize}
  \end{block}
  \bigskip
  \begin{block}{Cierre de la Unidad 2}
    Hemos cubierto el análisis exploratorio completo: estadística descriptiva univariada, distribuciones, detección de outliers, análisis bivariado y visualización en Python y R.
    La Unidad 3 abordará los métodos de \textbf{regresión}.
  \end{block}
\end{frame}

\end{document}
